\textbf{(i)} The ideal $\mathfrak{p}^nA_{\mathfrak{p}}=(\mathfrak{p}A_{\mathfrak{p}})^n$ is primary by Proposition 4.2. Its contraction is therefore $\mathfrak{p}$-primary.

\textbf{(ii)} Every associated prime of $\mathfrak{p}^n$ contains $r(\mathfrak{p}^n)=\mathfrak{p}$. Localization at $\mathfrak{p}$ discards all components except the unique $\mathfrak{p}$-primary component. Contracting recovers that component and gives $\mathfrak{p}^{(n)}$.

\textbf{(iii)} The radical of $\mathfrak{p}^{(m)}\mathfrak{p}^{(n)}$ is $\mathfrak{p}$, and
\[(\mathfrak{p}^{(m)}\mathfrak{p}^{(n)})A_{\mathfrak{p}}=(\mathfrak{p}A_{\mathfrak{p}})^{m+n}.\]
The argument in (ii) identifies its $\mathfrak{p}$-primary component with the contraction, namely $\mathfrak{p}^{(m+n)}$.

\textbf{(iv)} Equality implies primaryness by (i). Conversely, a primary ideal with radical $\mathfrak{p}$ equals the contraction of its extension to $A_{\mathfrak{p}}$.
